\begin{table*}[t]
\centering\small
\begin{tabular}{lcccccccc}
\toprule
 & \multicolumn{4}{c}{Books} & \multicolumn{4}{c}{DVDs} \\
\cmidrule(lr){2-5}\cmidrule(lr){6-9}
 & \multicolumn{2}{c}{Sparse LR} & \multicolumn{2}{c}{DT} & \multicolumn{2}{c}{Sparse LR} & \multicolumn{2}{c}{DT} \\
 & Ours & Paper & Ours & Paper & Ours & Paper & Ours & Paper \\
\midrule
Random & 17.1 & 17.4 & 14.6 & 20.6 & 13.9 & 19.2 & 13.1 & 17.4 \\
Parzen & 35.3 & 72.8 & 54.9 & 78.9 & 40.7 & 60.8 & 69.7 & 80.6 \\
Greedy & 45.5 & 64.3 & 28.0 & 37.0 & 51.2 & 63.4 & 38.1 & 47.6 \\
LIME & 86.3 & 92.1 & 89.7 & 97.0 & 88.3 & 90.2 & 95.6 & 97.8 \\
\midrule
Test acc. & 60.5 & -- & 63.0 & -- & 66.0 & -- & 64.8 & -- \\
\bottomrule
\end{tabular}

\caption{Replication of the faithfulness experiment (paper \S5.2, Figs.~6--7) on the authors' \texttt{processed\_acl} input: recall (\%) of the gold features by each explainer's top $K{=}10$, averaged over the 400 test reviews with a non-empty gold set. ``Paper'' values are the bar labels in the paper's figures. Test accuracy is that of the interpretable classifier, which the paper does not report.}
\label{tab:recallrepl}
\end{table*}

\begin{table*}[t]
\centering\small
\resizebox{\textwidth}{!}{\begin{tabular}{lcccccccc}
\toprule
 & \multicolumn{4}{c}{Books} & \multicolumn{4}{c}{DVDs} \\
\cmidrule(lr){2-5}\cmidrule(lr){6-9}
 & LR & NN & RF & SVM & LR & NN & RF & SVM \\
\midrule
Random & 14.7 {\scriptsize(14.6)} & 14.4 {\scriptsize(14.8)} & 14.7 {\scriptsize(14.7)} & 14.7 {\scriptsize(14.7)} & 14.9 {\scriptsize(14.2)} & 15.1 {\scriptsize(14.3)} & 14.8 {\scriptsize(14.5)} & 15.1 {\scriptsize(14.4)} \\
Parzen & 83.8 {\scriptsize(84.0)} & 90.4 {\scriptsize(87.6)} & 94.5 {\scriptsize(94.3)} & 89.5 {\scriptsize(92.3)} & 82.1 {\scriptsize(87.0)} & 80.4 {\scriptsize(81.7)} & 86.0 {\scriptsize(94.2)} & 82.4 {\scriptsize(87.3)} \\
Greedy & 54.5 {\scriptsize(53.7)} & 47.4 {\scriptsize(47.4)} & 48.3 {\scriptsize(45.0)} & 55.0 {\scriptsize(53.3)} & 53.6 {\scriptsize(52.4)} & 63.8 {\scriptsize(58.1)} & 60.4 {\scriptsize(46.6)} & 54.3 {\scriptsize(55.1)} \\
LIME & 96.5 {\scriptsize(96.6)} & 94.7 {\scriptsize(94.5)} & 98.3 {\scriptsize(96.2)} & 96.3 {\scriptsize(96.7)} & 96.7 {\scriptsize(96.6)} & 90.1 {\scriptsize(91.8)} & 98.6 {\scriptsize(96.1)} & 96.5 {\scriptsize(95.6)} \\
\midrule
Test acc. & 81.5 & 57.2 & 75.5 & 80.2 & 81.2 & 61.3 & 81.5 & 80.0 \\
Flipped & 11.3\% & 17.4\% & 9.7\% & 11.6\% & 11.4\% & 25.8\% & 10.7\% & 12.1\% \\
\bottomrule
\end{tabular}
}
\caption{Replication of the trust experiment (paper \S5.3, Table~1) on the authors' \texttt{processed\_acl} input: F1 (\%) of the trustworthy class, mean over 100 rounds with 25\% of the vocabulary untrustworthy. The paper's values are in brackets. ``Flipped'' is the share of test predictions the oracle marks untrustworthy. All differences between LIME and each baseline are significant ($t$-test, $p<10^{-4}$).}
\label{tab:trustrepl}
\end{table*}

\begin{table}[t]
\centering\small
\begin{tabular}{lcccc}
\toprule
 & \multicolumn{2}{c}{Trust F1 (8)} & \multicolumn{2}{c}{Recall (4)} \\
\cmidrule(lr){2-3}\cmidrule(lr){4-5}
Explainer & raw & acl & raw & acl \\
\midrule
Random & \MAERawTrustRandom & \MAEAclTrustRandom & \MAERawFaithRandom & \MAEAclFaithRandom \\
Parzen & \MAERawTrustParzen & \MAEAclTrustParzen & \MAERawFaithParzen & \MAEAclFaithParzen \\
Greedy & \MAERawTrustGreedy & \MAEAclTrustGreedy & \MAERawFaithGreedy & \MAEAclFaithGreedy \\
LIME & \MAERawTrustLime & \MAEAclTrustLime & \MAERawFaithLime & \MAEAclFaithLime \\
\bottomrule
\end{tabular}
\caption{Mean absolute difference from the paper (points) when the same code is run on our own tokenisation of the raw reviews (``raw'') and on the authors' pre-processed files (``acl''), over the 8 trust and 4 faithfulness settings.}
\label{tab:input}
\end{table}

\paragraph{Results.} Tables~\ref{tab:recallrepl} and~\ref{tab:trustrepl} give the replicated numbers next to the published ones. LIME is the best of the four explainers in all 12 settings (4 faithfulness, 8 trust), as the paper reports. The trust experiment replicates almost exactly. LIME's F1 is \ReplTrustRange{} against the paper's 91.8--96.7, a mean absolute difference of \ReplTrustMAE{} point, and all \ReplTrustWithinFive{} settings fall within the 5-point tolerance set in our proposal. Over all 32 explainer--setting pairs the mean difference is \ReplTrustAllMAE{} points, and several values are within 0.2 of the paper (e.g.\ Books LR: LIME 96.5 vs.\ 96.6, Parzen 83.8 vs.\ 84.0, Random 14.7 vs.\ 14.6). The ordering LIME $>$ Parzen $>$ Greedy $>$ Random holds in every trust setting, and per-round standard deviations are about 1 point (Table~\ref{tab:prf}). In the faithfulness experiment, LIME recovers \ReplRecallRange{}\% of the gold features against the paper's 90.2--97.8\%, \ReplRecallGap{} points lower on average. The paper's qualitative findings hold here too: Random is near chance, Greedy is much weaker on trees than on LR (and matches the paper's tree values, 37.1 vs.\ 37.0 on Books), and LIME leads every baseline.

\paragraph{The input matters more than the code.} Running the same code on our own tokenisation of the raw reviews (\S\ref{sec:data}) gives the same conclusions, but is consistently further from the paper (Table~\ref{tab:input}). The effect is small for LIME (\MAERawTrustLime{} vs.\ \MAEAclTrustLime{} points in trust), whose kernel only looks at the document's own words. It is large for Parzen's faithfulness (\MAERawFaithParzen{} vs.\ \MAEAclFaithParzen{} points). Parzen uses a squared-Euclidean kernel over the whole $\sim$20k-word vocabulary with bandwidths the authors tuned on their own files, so on raw text (about 110 instead of 85 distinct words per review) those bandwidths are wrong. Our bandwidth ablation on raw text (\texttt{src/ablation\_parzen\_sigma.py}) agrees: even the best $\sigma$, chosen on the test set, recovers only part of the gap (51--71\% recall), while the authors' input brings Parzen to within \MAEAclFaithParzen{} points of the paper on average.

\paragraph{Remaining differences.} \emph{(1) LIME's recall for the Books sparse LR} is 85.1 against 92.1 (95\% CI $\pm$2.5). With $C{=}0.03$ the model keeps only 6 non-zero words (test accuracy 58\%), so each document has about two gold words and one miss costs 50\%. Which 6 words are selected may depend on solver details that changed between scikit-learn 0.18 and 1.9. We consider this the most likely cause, but could not test it without the original stack. \emph{(2) The decision trees are unpruned (depth 49--53)}, as in the released code, although the paper says no prediction uses more than 10 features and its figure files are named \texttt{shallow\_tree}. We tested a depth-10 tree (\texttt{--tree-max-depth 10}). It changes LIME's recall by only about a point, in opposite directions (Books \BooksDeepTreeLime{}$\rightarrow$\BooksShallowTreeLime{}, DVDs \DvdDeepTreeLime{}$\rightarrow$\DvdShallowTreeLime{}), and it moves Greedy far from the paper (Books 60.2 vs.\ 37.0, DVDs 58.9 vs.\ 47.6). So the paper most likely used the unpruned tree, and tree depth is ruled out as a cause. \emph{(3) Parzen's trust F1 on DVDs} is 1--8 points low, most of all with RF (86.0 vs.\ 94.2). The authors' DVD bandwidths are small ($\sigma{=}0.5$--$0.75$, against 0.5--7 for Books), so the Parzen estimate rests on very few neighbours and is sensitive to small differences in the data; we did not test this further. \emph{(4) RF} uses bag-of-words features, as in the released code, instead of the embedding forest described in the paper. This is the most likely reason its Greedy values differ most (60.4 vs.\ 46.6 on DVDs).
